\documentclass[11pt,a4paper]{article}

% -------------------------------------------------------
% PACCHETTI
% -------------------------------------------------------
% babel italian non disponibile in questa installazione
% \usepackage[italian]{babel}
\usepackage[utf8]{inputenc}
\usepackage[T1]{fontenc}
\usepackage{amsmath, amssymb, amsthm}
\usepackage{mathtools}
\usepackage{geometry}
\usepackage{hyperref}
\usepackage{xcolor}
\usepackage{tcolorbox}
% algoritmi implementati manualmente con tcolorbox + enumerate
\usepackage{booktabs}
\usepackage{enumitem}
% \usepackage{microtype}  % richiede font scalabili

\geometry{margin=2.5cm}
\hypersetup{colorlinks=true, linkcolor=blue!70!black, urlcolor=blue!70!black}

% -------------------------------------------------------
% STILI TEOREMI
% -------------------------------------------------------
\theoremstyle{definition}
\newtheorem{definition}{Definizione}[section]
\newtheorem{example}{Esempio}[section]

\theoremstyle{plain}
\newtheorem{theorem}{Teorema}[section]
\newtheorem{proposition}{Proposizione}[section]
\newtheorem{lemma}{Lemma}[section]
\newtheorem{corollary}{Corollario}[section]

\theoremstyle{remark}
\newtheorem{remark}{Osservazione}[section]

% -------------------------------------------------------
% SCATOLE COLORATE
% -------------------------------------------------------
\tcbuselibrary{breakable, skins}

\newtcolorbox{keybox}[1][]{
  colback=blue!5!white,
  colframe=blue!60!black,
  fonttitle=\bfseries,
  title={#1},
  breakable
}

\newtcolorbox{algobox}[1][]{
  colback=green!5!white,
  colframe=green!50!black,
  fonttitle=\bfseries,
  title={#1},
  breakable
}

\newtcolorbox{intubox}[1][Intuizione]{
  colback=orange!5!white,
  colframe=orange!60!black,
  fonttitle=\bfseries,
  title={#1},
  breakable
}

% -------------------------------------------------------
% AMBIENTE PER PSEUDOCODICE (senza pacchetto algorithm)
% -------------------------------------------------------
\newenvironment{pseudocode}{%
  \begin{list}{}{%
    \setlength{\leftmargin}{1.5em}%
    \setlength{\itemsep}{2pt}%
    \setlength{\parsep}{0pt}%
    \setlength{\topsep}{4pt}%
  }%
}{%
  \end{list}%
}
\newcommand{\pcline}[1]{\item[] #1}
\newcommand{\pcindent}{\hspace{1.5em}}
\newcommand{\pcindentII}{\hspace{3em}}

% -------------------------------------------------------
% MACRO
% -------------------------------------------------------
\newcommand{\R}{\mathbb{R}}
\newcommand{\norm}[1]{\left\| #1 \right\|}
\newcommand{\abs}[1]{\left| #1 \right|}
\newcommand{\T}{^{\top}}
\newcommand{\grad}{\nabla}
\newcommand{\subgr}{\partial}
\newcommand{\argmin}{\operatorname*{argmin}}

% -------------------------------------------------------
% DOCUMENTO
% -------------------------------------------------------
\begin{document}

% -------------------------------------------------------
% FRONTESPIZIO
% -------------------------------------------------------
\begin{titlepage}
  \centering
  \vspace*{3cm}
  {\LARGE\bfseries Manuale Teorico\par}
  \vspace{0.5cm}
  {\Large Progetto 25 --- ELM con Regolarizzazione $\ell_1$\par}
  \vspace{1cm}
  {\large Corso: Computational Mathematics for Learning and Data Analysis\par}
  \vspace{0.5cm}
  {\large A.A.\ 2025/2026\par}
  \vspace{3cm}
  \begin{abstract}
    Questo manuale raccoglie gli argomenti teorici necessari per lo sviluppo e la comprensione
    del \textbf{Progetto 25} del corso CM.
    Il progetto riguarda il training di una \emph{Extreme Learning Machine} (ELM)
    con regolarizzazione $\ell_1$ (LASSO), risolto tramite due algoritmi:
    \textbf{A1 = IRLS} (Iteratively Reweighted Least Squares) e
    \textbf{A2 = Metodo del Sottogradiente Deflesso con Target Level}.
    Il documento è strutturato come percorso di studio progressivo: si parte dal
    problema di minimi quadrati lineari, si analizzano i solutori numerici
    (fattorizzazione di Cholesky e gradiente coniugato), si introducono convessità e
    ottimalità per funzioni non differenziabili, e infine si descrivono i due algoritmi.
  \end{abstract}
  \vfill
  {\small Basato sulle lezioni del corso:
  \emph{Optimization} (OPT/3, OPT/4, OPT/5) e
  \emph{Numerical Linear Algebra} (NLA/3, NLA/4, NLA/5, NLA/20).}
\end{titlepage}

\tableofcontents
\newpage

% ================================================================
\section{Extreme Learning Machine e Formulazione del Problema}
\label{sec:elm}
% ================================================================

\subsection{Architettura ELM}

Una \textbf{Extreme Learning Machine} (ELM) è una rete neurale con un singolo strato nascosto
in cui i pesi del primo strato $W_1$ (e i bias) sono \emph{fissi e generati casualmente},
mentre solo i pesi di output $w \in \R^n$ vengono ottimizzati.

Dato un dataset di addestramento $\{(x_i, y_i)\}_{i=1}^m$ con $x_i \in \R^d$ e $y_i \in \R$,
la rete calcola:
\[
  \hat{y}_i = \sum_{j=1}^n \sigma(W_{1,j}\, x_i + b_j)\, w_j,
\]
dove $\sigma$ è una funzione di attivazione non lineare (es.\ sigmoide o ReLU).

\begin{keybox}[Rappresentazione matriciale]
  Si definisce la \textbf{matrice delle feature} $X \in \R^{m \times n}$:
  \[
    X_{ij} = \sigma(W_{1,j}\, x_i + b_j).
  \]
  Poiché $W_1$ è fissato, $X$ è nota.
  Il vettore target è $y \in \R^m$, e l'output della rete si scrive compattamente come $\hat{y} = Xw$.
  Il training dell'ELM si riduce a un problema di \textbf{ottimizzazione lineare} nei pesi $w$.
\end{keybox}

\subsection{Problema di Ottimizzazione con Regolarizzazione $\ell_1$}

Addestrare l'ELM minimizzando l'errore quadratico con regolarizzazione LASSO equivale a risolvere:

\begin{keybox}[Problema Centrale del Progetto]
  \[
    \min_{w \in \R^n}\; f(w)
    \;=\; \underbrace{\norm{Xw - y}_2^2}_{\text{errore quadratico}}
    \;+\; \lambda\, \underbrace{\norm{w}_1}_{\text{regolarizzazione }\ell_1},
    \qquad \lambda > 0.
  \]
\end{keybox}

\begin{intubox}
  La norma $\ell_1$ \textbf{promuove la sparsità}: spinge molti pesi $w_i$ ad essere
  \emph{esattamente zero}, selezionando automaticamente le feature più rilevanti.
  Questo è particolarmente utile quando $n \gg m$ (molte feature, pochi campioni),
  scenario comune nelle ELM ad alto numero di neuroni nascosti.
  In contrasto, la regolarizzazione $\ell_2$ (Ridge) distribuisce uniformemente i pesi senza azzerarli.
\end{intubox}

\begin{remark}
  La funzione obiettivo $f(w)$ è la somma di una parte \emph{liscia} (il termine quadratico,
  differenziabile ovunque) e di una parte \emph{non liscia} ($\lambda\norm{w}_1$, non
  differenziabile nei punti con $w_i = 0$).
  Questa struttura mista richiede strumenti distinti: solutori lineari diretti/iterativi
  per il termine quadratico, e tecniche di ottimizzazione non liscia per la norma $\ell_1$.
\end{remark}

% ================================================================
\section{Minimi Quadrati Lineari}
\label{sec:ls}
% ================================================================

\subsection{Formulazione e Geometria}

Dato $A \in \R^{m \times n}$ (con $m \geq n$) e $y \in \R^m$, il \textbf{problema dei minimi quadrati} è:
\[
  \min_{w \in \R^n} \norm{Aw - y}_2.
\]

\begin{intubox}
  Il vettore $Aw$ vive nello spazio colonna $\operatorname{Im}(A)$.
  Se $y \notin \operatorname{Im}(A)$, non esiste soluzione esatta.
  Il problema chiede il punto in $\operatorname{Im}(A)$ più vicino a $y$:
  tale punto è la \textbf{proiezione ortogonale} di $y$ su $\operatorname{Im}(A)$.
  Il vettore residuo $r = Aw_* - y$ è ortogonale a ogni colonna di $A$.
\end{intubox}

\subsection{Rango Pieno per Colonne e Unicità}

\begin{definition}[Rango Pieno per Colonne]
  $A \in \R^{m \times n}$ ha \textbf{rango pieno per colonne} se $\ker(A) = \{0\}$,
  ovvero $\operatorname{rk}(A) = n$.
\end{definition}

\begin{theorem}
  Il problema dei minimi quadrati ha soluzione \textbf{unica} se e solo se $A$ ha rango pieno per
  colonne. Inoltre, $A$ ha rango pieno per colonne se e solo se $A\T A$ è \textbf{definita positiva}.
\end{theorem}

\begin{proof}
  Per ogni $z \neq 0$: $z\T A\T A z = \norm{Az}^2 > 0 \iff Az \neq 0 \iff z \notin \ker(A)$.
\end{proof}

\subsection{Equazioni Normali}

Minimizzare $\norm{Aw-y}^2$ equivale a minimizzare:
\[
  f(w) = \tfrac{1}{2}w\T A\T A\, w - (A\T y)\T w + \tfrac{1}{2}\norm{y}^2.
\]
Il gradiente è $\grad f(w) = A\T A w - A\T y$ e l'Hessiano è $A\T A \succ 0$ (funzione
strettamente convessa). Il minimo si trova ponendo il gradiente a zero:

\begin{keybox}[Equazioni Normali]
  \[
    A\T A\, w_* = A\T y.
  \]
  La matrice $Q = A\T A$ è \textbf{simmetrica definita positiva} (SPD).
  La soluzione è $w_* = (A\T A)^{-1} A\T y =: A^+ y$,
  dove $A^+ = (A\T A)^{-1} A\T$ è la \textbf{pseudoinversa} di Moore-Penrose.
\end{keybox}

\begin{remark}
  La denominazione ``normali'' viene dal fatto che il residuo $Aw_* - y$ è \emph{normale}
  (perpendicolare) allo spazio $\operatorname{Im}(A)$:
  $(Av)\T(Aw_* - y) = 0$ per ogni $v \in \R^n$, che riscritta dà $A\T(Aw_* - y) = 0$.
\end{remark}

\subsection{Costo Computazionale}

Per $A \in \R^{m \times n}$ con $m > n$, sfruttando la simmetria di $A\T A$:
\begin{enumerate}
  \item Calcolo di $A\T A$: costo $mn^2$ operazioni.
  \item Calcolo di $A\T y$: costo $2mn$ (trascurabile).
  \item Risoluzione di $Qw = b$ con Cholesky: costo $\frac{1}{3}n^3$.
\end{enumerate}
Il costo totale è $O(mn^2)$ per $m \gg n$.

% ================================================================
\section{Fattorizzazione di Cholesky}
\label{sec:chol}
% ================================================================

\subsection{Sfruttare la Simmetria: da LU a $LDL\T$}

La fattorizzazione LU per matrici generali $n \times n$ ha costo $\frac{2}{3}n^3 + O(n^2)$.
Per matrici \textbf{simmetriche} si può fare molto meglio applicando l'eliminazione in modo simmetrico,
producendo la fattorizzazione $LDL\T$.

\begin{keybox}[Fattorizzazione $LDL\T$]
  Per ogni matrice simmetrica $M = M\T \in \R^{n \times n}$ (per cui l'algoritmo non incontra pivot nulli):
  \[
    M = L D L\T,
  \]
  dove $L$ è triangolare inferiore con $1$ sulla diagonale e $D$ è diagonale.
  \textbf{Costo:} $\frac{1}{3}n^3 + O(n^2)$ --- la \emph{metà} di LU.
\end{keybox}

\subsection{Cholesky per Matrici SPD}

Per matrici SPD tutti i pivot sono positivi ($D_{kk} > 0$), quindi $D^{1/2} = \operatorname{diag}(D_{11}^{1/2}, \ldots, D_{nn}^{1/2})$
è ben definita. Si ottiene:

\begin{keybox}[Fattorizzazione di Cholesky]
  Per ogni matrice SPD $M \in \R^{n \times n}$:
  \[
    M = L D L\T = L D^{1/2}(D^{1/2} L\T) = R\T R,
  \]
  dove $R = D^{1/2} L\T$ è triangolare superiore con elementi diagonali positivi.
  In MATLAB: \texttt{R = chol(A)}, che restituisce $R$ tale che $A = R\T R$.
\end{keybox}

\begin{proposition}[Stabilità]
  Si può dimostrare (tramite la SVD di $R$) che $\norm{R} = \norm{M}^{1/2}$.
  La fattorizzazione di Cholesky è \textbf{backward stable} per tutte le matrici SPD
  e non richiede pivoting.
\end{proposition}

\subsubsection*{Risoluzione di $Mw = b$ con Cholesky}

Avendo $M = R\T R$, il sistema si riduce a due triangolari (forward/backward substitution), ognuna $O(n^2)$:
\[
  R\T z = b \;\Rightarrow\; z \quad\text{(forward)}, \qquad Rw = z \;\Rightarrow\; w \quad\text{(backward)}.
\]

\begin{intubox}
  Cholesky è il metodo ottimale per sistemi SPD: stabile senza pivoting, costo $\frac{1}{3}n^3$
  (metà di LU), ideale per le equazioni normali $A\T Aw = A\T y$.
  Per la regola pratica: ``\emph{Is your matrix posdef? Use Cholesky}.''
\end{intubox}

% ================================================================
\section{Gradiente Coniugato (CG)}
\label{sec:cg}
% ================================================================

\subsection{Motivazione: Sistemi Sparsi di Grandi Dimensioni}

Per sistemi SPD di grandi dimensioni e \textbf{sparsi} (pochi elementi non nulli), i metodi diretti
come Cholesky hanno costo $O(n^3)$ e memoria $O(n^2)$: inapplicabili.
Il \textbf{Gradiente Coniugato} risolve $Qx = v$ (con $Q \succ 0$) iterativamente,
usando \emph{un solo prodotto matrice-vettore per iterazione}.

\subsection{Riformulazione come Ottimizzazione}

Risolvere $Qx = v$ ($Q \succ 0$) è equivalente a minimizzare:
\[
  f(x) = \tfrac{1}{2} x\T Q x - v\T x + \text{const}.
\]

\subsection{Spazi di Krylov e Ottimalità nel Sottospazio}

\begin{definition}[Spazio di Krylov]
  Dati $Q \in \R^{m \times m}$ e $v \in \R^m$, lo \textbf{spazio di Krylov} di ordine $k$ è:
  \[
    K_k(Q, v) = \operatorname{span}\{v,\, Qv,\, Q^2v,\, \ldots,\, Q^{k-1}v\}.
  \]
\end{definition}

Il CG computa iterati con la \textbf{proprietà di ottimalità nel sottospazio}:
$x_k = \argmin_{x \in K_k(Q,v)} f(x)$.

\subsection{Algoritmo CG e $Q$-Ortogonalità}

\begin{algobox}[Algoritmo: Gradiente Coniugato]
  \begin{pseudocode}
    \pcline{$x_0 \leftarrow 0$,\quad $r_0 \leftarrow d_0 \leftarrow v$}
    \pcline{\textbf{for} $j = 1, 2, \ldots, m$:}
    \pcline{\pcindent $\alpha_j \leftarrow (r_{j-1}\T r_{j-1}) / (d_{j-1}\T Q\, d_{j-1})$
      \hfill \textit{ricerca lineare esatta}}
    \pcline{\pcindent $x_j \leftarrow x_{j-1} + \alpha_j\, d_{j-1}$}
    \pcline{\pcindent $r_j \leftarrow r_{j-1} - \alpha_j\, Q\, d_{j-1}$
      \hfill \textit{aggiornamento residuo}}
    \pcline{\pcindent $\beta_j \leftarrow (r_j\T r_j) / (r_{j-1}\T r_{j-1})$}
    \pcline{\pcindent $d_j \leftarrow r_j + \beta_j\, d_{j-1}$
      \hfill \textit{direzione $Q$-ortogonale}}
    \pcline{\textbf{end for}}
  \end{pseudocode}
\end{algobox}

\textbf{Proprietà chiave:}
\begin{itemize}[noitemsep]
  \item I residui sono ortogonali: $r_i\T r_j = 0$ per $i \neq j$.
  \item Le direzioni sono $Q$-ortogonali: $d_i\T Q\, d_j = 0$ per $i \neq j$.
  \item Convergenza garantita in $\leq m$ iterazioni (aritmetica esatta).
  \item Costo per iterazione: $O(\text{nnz}(Q))$, un solo prodotto matrice-vettore.
  \item Storage: solo 3 vettori ($x_j$, $r_j$, $d_j$).
\end{itemize}

\subsection{Convergenza e Numero di Condizionamento}

\begin{definition}[Numero di Condizionamento]
  Per $Q \succ 0$:
  \[
    \kappa(Q) = \frac{\lambda_{\max}(Q)}{\lambda_{\min}(Q)} \geq 1.
  \]
\end{definition}

\begin{theorem}[Convergenza del CG]
  Dopo $k$ iterazioni:
  \[
    \norm{x_k - x_*}_{Q} \leq 2\left(\frac{\sqrt{\kappa}-1}{\sqrt{\kappa}+1}\right)^k \norm{x_0 - x_*}_Q,
  \]
  dove $\norm{z}_Q = (z\T Qz)^{1/2}$ e $\kappa = \kappa(Q)$.
\end{theorem}

\begin{intubox}
  Per ottenere $\varepsilon$-accuratezza nella norma $Q$: bastano $O(\sqrt{\kappa}\, \log(1/\varepsilon))$ iterazioni.
  \textbf{Attenzione:} per le equazioni normali $\kappa(A\T A) = \kappa(A)^2$, quindi il sistema normale
  è molto più mal condizionato della matrice originale.
  Il precondizionamento serve a ridurre il $\kappa$ effettivo.
\end{intubox}

% ================================================================
\section{Convessità e Condizioni di Ottimalità}
\label{sec:opt}
% ================================================================

\subsection{Definizioni di Convessità}

\begin{definition}[Convessità]
  $f : \R^n \to \R$ è \textbf{convessa} se per ogni $x, z \in \R^n$ e $\alpha \in [0,1]$:
  \[
    f(\alpha x + (1-\alpha)z) \leq \alpha f(x) + (1-\alpha)f(z).
  \]
  Equivalentemente (per $f \in C^1$): $f(z) \geq f(x) + \langle\grad f(x), z-x\rangle$.
  Equivalentemente (per $f \in C^2$): $\grad^2 f(x) \succeq 0$.
\end{definition}

\begin{definition}[Forte Convessità]
  $f$ è \textbf{$\tau$-fortemente convessa} ($\tau > 0$) se:
  \[
    f(z) \geq f(x) + \langle\grad f(x), z-x\rangle + \frac{\tau}{2}\norm{z-x}^2.
  \]
  Per $f \in C^2$: equivale a $\grad^2 f(x) \succeq \tau I$, ovvero $\lambda_{\min}(\grad^2 f) \geq \tau > 0$.
\end{definition}

\begin{definition}[$L$-Lisciezza]
  $f \in C^1$ è \textbf{$L$-liscia} se il gradiente è $L$-Lipschitz:
  \[
    \norm{\grad f(x) - \grad f(z)} \leq L\norm{x-z}, \quad \forall x, z.
  \]
  Per $f \in C^2$: equivale a $\grad^2 f(x) \preceq LI$, ovvero $\lambda_{\max}(\grad^2 f) \leq L$.
\end{definition}

\begin{keybox}[Minimi quadrati come funzione SPD]
  Per $f(w) = \frac{1}{2}\norm{Xw-y}^2$, l'Hessiano è $\grad^2 f = X\T X$.
  Quindi $f$ è $L$-liscia con $L = \lambda_{\max}(X\T X)$ e $\tau$-fortemente convessa con
  $\tau = \lambda_{\min}(X\T X)$. Il numero di condizionamento è $\kappa = L/\tau = \kappa(X\T X)$.
\end{keybox}

\subsection{Convergenza del Metodo del Gradiente con Forte Convessità}

Per $f$ $L$-liscia e $\tau$-fortemente convessa, il metodo del gradiente con passo fisso
$\alpha = 2/(L+\tau)$ converge \textbf{linearmente}:
\[
  \norm{x^{k+1} - x_*} \leq r^k \norm{x^1 - x_*},
  \qquad r = \frac{\kappa - 1}{\kappa + 1} < 1, \quad \kappa = L/\tau.
\]

\subsection{Condizioni di Ottimalità per Funzioni Lisce}

Per $f \in C^1$ convessa: $w^*$ è minimo globale $\iff \grad f(w^*) = 0$.

\subsection{Sottogradiente e Subdifferenziale}

Per funzioni \textbf{non differenziabili} (come $\norm{w}_1$), il gradiente non esiste in tutti i punti.
Si introduce la seguente generalizzazione.

\begin{definition}[Sottogradiente]
  $g \in \R^n$ è un \textbf{sottogradiente} di $f$ in $x$ se:
  \[
    f(z) \geq f(x) + \langle g, z-x\rangle, \quad \forall z \in \R^n.
  \]
  A differenza del gradiente, un sottogradiente \emph{non} garantisce la direzione di discesa:
  può accadere che $f(x - \alpha g) \geq f(x)$ per ogni $\alpha > 0$.
\end{definition}

\begin{definition}[Subdifferenziale]
  Il \textbf{subdifferenziale} di $f$ in $x$ è l'insieme di tutti i sottogradienti:
  \[
    \subgr f(x) = \bigl\{ g \in \R^n : f(z) \geq f(x) + \langle g, z-x\rangle \;\forall z \bigr\}.
  \]
\end{definition}

\begin{example}[Valore assoluto e norma $\ell_1$]
  Per $f(w) = \norm{w}_1$:
  \[
    [\subgr f(w)]_i =
    \begin{cases}
      \{+1\}  & \text{se } w_i > 0, \\
      \{-1\}  & \text{se } w_i < 0, \\
      [-1,+1] & \text{se } w_i = 0.
    \end{cases}
  \]
  Dove $f$ è differenziabile ($w_i \neq 0$), il subdifferenziale è un singoletto (il gradiente);
  dove non lo è ($w_i = 0$), è un intervallo chiuso.
\end{example}

\begin{keybox}[Condizione di Ottimalità Generale]
  Per $f$ convessa (differenziabile o no): $w^*$ è minimo globale se e solo se:
  \[
    0 \in \subgr f(w^*).
  \]
\end{keybox}

\subsubsection*{Regola di Somma per il Subdifferenziale}

Se $f = f_1 + f_2$ con $f_1$ differenziabile e $f_2$ convessa:
\[
  \subgr f(w) = \grad f_1(w) + \subgr f_2(w).
\]
Per la funzione obiettivo del progetto $f(w) = \norm{Xw-y}^2 + \lambda\norm{w}_1$:
\[
  \subgr f(w) = 2X\T(Xw - y) + \lambda\, \subgr\norm{w}_1.
\]

% ================================================================
\section{Ottimizzazione Non Liscia}
\label{sec:nonsmooth}
% ================================================================

\subsection{Perché i Metodi Lisci Falliscono}

Per $f \notin C^1$ come $\norm{w}_1$, il metodo del gradiente \textbf{non funziona}:
\begin{enumerate}
  \item Il passo fisso non converge: per $f(x) = L|x|$ con $x^0 = -\alpha L/2$,
    iterando si ottiene un ciclo $x^0 \to x^1 \to x^0 \to \cdots$ senza convergenza.
  \item La norma del sottogradiente non è un indicatore di ottimalità:
    $\norm{g}$ piccola non implica $f(x) \approx f_*$, né viceversa.
    Non c'è un criterio di arresto affidabile.
\end{enumerate}

\begin{keybox}[Complessità dell'Ottimizzazione Non Liscia]
  \begin{center}
    \begin{tabular}{llll}
      \toprule
      Classe & Convessità & Regolarità & Complessità \\
      \midrule
      $f \in C^1$ & $\tau$-convessa & $L$-liscia & $O(\log(1/\varepsilon))$ \\
      $f \notin C^1$ & $\tau$-convessa & $L$-Lipschitz & $\Omega(L^2/\varepsilon^2)$ \\
      $f \in C^1$ & convessa & $L$-liscia & $O(1/\sqrt{\varepsilon})$ \\
      $f \notin C^1$ & convessa & $L$-Lipschitz & $\Omega(L/\varepsilon^2)$ \\
      \bottomrule
    \end{tabular}
  \end{center}
  \vspace{0.3em}
  L'ottimizzazione non liscia è \textbf{ordini di grandezza più lenta} di quella liscia.
\end{keybox}

\subsection{Metodo del Sottogradiente con Passo DSS}

L'iterazione base del metodo del sottogradiente è:
\[
  x^{i+1} = x^i - \alpha^i g^i, \qquad g^i \in \subgr f(x^i).
\]

\begin{theorem}[Relazione Fondamentale]
  Per ogni $g^i \in \subgr f(x^i)$:
  \[
    \norm{x^{i+1} - x_*}^2
    \leq \norm{x^i - x_*}^2 + 2\alpha^i\bigl(f_* - f(x^i)\bigr) + (\alpha^i)^2\norm{g^i}^2.
  \]
\end{theorem}

Per $\alpha^i$ piccolo, il termine negativo $2\alpha^i(f_* - f(x^i))$ domina e il passo avvicina a $x_*$.
Il passo non può essere fisso (diverge); deve tendere a zero, ma non troppo in fretta.

\begin{definition}[Passo DSS (Diminishing-Square-Summable)]
  \[
    \sum_{i=1}^\infty \alpha^i = \infty \quad \text{e} \quad
    \sum_{i=1}^\infty (\alpha^i)^2 < \infty.
  \]
  Esempio: $\alpha^i = 1/i$.
\end{definition}

Con il passo DSS: $\exists k$ tale che $f(x^k) - f_* \leq \varepsilon$ (convergenza sulle \emph{iterate record}).
Tuttavia la convergenza pratica è molto lenta.

\subsection{Passo di Polyak}

Se si conosce $f_*$, si può ottimizzare il passo derivando dalla relazione fondamentale:
\[
  \alpha^i_{\text{opt}} = \argmin_{\alpha \geq 0}
    \bigl[\norm{x^i - x_*}^2 + 2\alpha(f_* - f^i) + \alpha^2\norm{g^i}^2\bigr]
  = \frac{f(x^i) - f_*}{\norm{g^i}^2}.
\]

\begin{keybox}[Passo di Polyak (PSS)]
  \[
    \alpha^i = \frac{f(x^i) - f_*}{\norm{g^i}^2} \geq 0.
  \]
  Con questo passo: $\norm{x^{i+1} - x_*}^2 < \norm{x^i - x_*}^2$ (monotono decrescente).

  Complessità: $O(1/\varepsilon^2)$ sulle iterate record, con costante molto migliore di DSS.
\end{keybox}

\begin{intubox}
  Il passo di Polyak è ``ottimale'' ma richiede la conoscenza di $f_*$, che in pratica non si ha.
  L'approccio \emph{target level} stima $f_*$ adattativamente, rifinanando la stima se troppo
  ottimistica.
\end{intubox}

% ================================================================
\section{Algoritmo A2: Sottogradiente Deflesso con Target Level}
\label{sec:a2}
% ================================================================

\subsection{Deflettere la Direzione}

Il metodo del sottogradiente puro usa $d^i = -g^i$.
Il metodo \textbf{deflesso} combina il sottogradiente corrente con la direzione precedente,
analogamente al gradiente coniugato:
\[
  d^i = \gamma^i g^i + (1 - \gamma^i) d^{i-1}, \qquad \gamma^i \in [0, 1].
\]
Il passo di aggiornamento è $x^{i+1} = x^i - \alpha^i d^i$.

\begin{intubox}
  Si tratta di un ``sottogradiente coniugato'': sfrutta la memoria delle iterazioni precedenti
  per costruire una direzione più informata. La direzione deflessa approssima una
  direzione di discesa migliore, in modo analogo al gradiente coniugato per funzioni quadratiche.
\end{intubox}

\subsection{Scelta Ottimale di $\gamma^i$}

Il parametro $\gamma^i$ viene scelto minimizzando la norma quadratica della direzione deflessa:

\begin{keybox}[Formula Chiusa per $\gamma^i$]
  \[
    \gamma^i = \argmin_{\gamma \in [0,1]} \norm{\gamma g^i + (1-\gamma) d^{i-1}}^2.
  \]
  Svolgendo il quadrato, si ottiene un problema quadratico univariato in $\gamma$ con soluzione:
  \[
    \gamma^i_{\text{unc}}
    = \frac{\langle d^{i-1}, d^{i-1} - g^i \rangle}{\norm{g^i - d^{i-1}}^2}
    = \frac{\norm{d^{i-1}}^2 - \langle g^i, d^{i-1}\rangle}{\norm{g^i - d^{i-1}}^2},
  \]
  e $\gamma^i = \operatorname{clip}(\gamma^i_{\text{unc}}, 0, 1)$.
\end{keybox}

\subsection{Meccanismo Target Level}

Il \emph{target level} stima adattativamente $f_*$ per il passo di Polyak.
Si mantiene un \textbf{valore di riferimento} $f_{\text{ref}}$ e una \textbf{soglia} $\delta > 0$.

\begin{keybox}[Logica del Target Level]
  \begin{itemize}[noitemsep]
    \item Il \textbf{target level} è $f_{\text{ref}} - \delta$ (livello da raggiungere).
    \item Il passo di Polyak usa il target al posto di $f_*$:
      $\alpha^i = \beta^i (f(x^i) - (f_{\text{ref}} - \delta)) / \norm{d^i}^2$.
    \item \textbf{Miglioramento buono} ($f(x) \leq f_{\text{ref}} - \delta/2$):
      si abbassa $f_{\text{ref}}$ e si azzera il contatore $r$.
    \item \textbf{Troppi passi senza miglioramento} ($r > R$):
      il target era troppo ambizioso; si riduce $\delta \leftarrow \rho\delta$ (target più facile).
  \end{itemize}
\end{keybox}

\subsection{Algoritmo Completo (A2)}

\begin{algobox}[Algoritmo A2: Sottogradiente Deflesso con Target Level]
  \textbf{Input:} $f$, $x^0$, $i_{\max}$, $\beta \in (0,2)$, $\delta_0 > 0$, $R > 0$, $\rho \in (0,1)$\\
  \textbf{Output:} Valore ottimale approssimato $\bar{f}$
  \begin{pseudocode}
    \pcline{$r \leftarrow 0$;\; $\delta \leftarrow \delta_0$;\;
      $f_{\text{ref}} \leftarrow \bar{f} \leftarrow f(x^0)$;\; $i \leftarrow 1$}
    \pcline{\textbf{while} $i < i_{\max}$:}
    \pcline{\pcindent $g^i \in \subgr f(x^i)$
      \hfill \textit{calcola un sottogradiente}}
    \pcline{\pcindent $\gamma^i \leftarrow \argmin_{\gamma \in [0,1]}\norm{\gamma g^i + (1-\gamma)d^{i-1}}^2$
      \hfill \textit{coefficiente di deflazione}}
    \pcline{\pcindent $d^i \leftarrow \gamma^i g^i + (1-\gamma^i)d^{i-1}$
      \hfill \textit{direzione deflessa}}
    \pcline{\pcindent $\alpha^i \leftarrow \beta^i \cdot \dfrac{f(x^i) - (f_{\text{ref}} - \delta)}{\norm{d^i}^2}$
      \hfill \textit{passo Polyak con target}}
    \pcline{\pcindent $x \leftarrow x^i - \alpha^i d^i$}
    \pcline{\pcindent \textbf{if} $f(x) \leq f_{\text{ref}} - \delta/2$:
      \hfill \textit{buon miglioramento}}
    \pcline{\pcindentII $f_{\text{ref}} \leftarrow \bar{f}$;\; $r \leftarrow 0$}
    \pcline{\pcindent \textbf{else if} $r > R$:
      \hfill \textit{target troppo ambizioso: abbassa soglia}}
    \pcline{\pcindentII $\delta \leftarrow \rho\delta$;\; $r \leftarrow 0$}
    \pcline{\pcindent \textbf{else}:}
    \pcline{\pcindentII $r \leftarrow r + \alpha^i\norm{d^i}$
      \hfill \textit{accumula progresso}}
    \pcline{\pcindent \textbf{end if}}
    \pcline{\pcindent $\bar{f} \leftarrow \min(\bar{f},\, f(x))$;\; $i \leftarrow i + 1$}
    \pcline{\textbf{end while}}
  \end{pseudocode}
\end{algobox}

\subsection{Note sulla Convergenza}

\begin{proposition}
  Sotto condizioni appropriate sui parametri ($\beta^i$, $\gamma^i$ legati dalla condizione
  ``stepsize-restricted'' o ``deflection-restricted''), la sequenza delle iterate record
  $\bar{f}^i$ converge a $f_*$ con complessità $O(1/\varepsilon^2)$.
\end{proposition}

Il calcolo del sottogradiente $g^i \in \subgr f(w^i)$ per $f(w) = \norm{Xw-y}^2 + \lambda\norm{w}_1$ è:
\[
  g^i = 2X\T(Xw^i - y) + \lambda\, s^i,
  \qquad s^i_j = \begin{cases} +1 & w^i_j > 0 \\ -1 & w^i_j < 0 \\ \in [-1,+1] & w^i_j = 0 \end{cases}
\]

% ================================================================
\section{Algoritmo A1: IRLS (Iteratively Reweighted Least Squares)}
\label{sec:a1}
% ================================================================

\subsection{Motivazione: Approssimare $\ell_1$ con Minimi Quadrati Pesati}

La norma $\norm{w}_1 = \sum_i |w_i|$ non è differenziabile. L'idea fondamentale dell'IRLS
è di approssimare $|w_i|$ usando il valore dell'iterata precedente $w_k$:
\[
  |w_i| = \frac{w_i^2}{|w_i|} \approx \frac{w_i^2}{\max(|w_{k,i}|,\,\varepsilon_{\text{thr}})}.
\]
In questo modo:
\[
  \norm{w}_1 \approx \sum_i \frac{w_i^2}{|w_{k,i}|} = w\T D_k w = \norm{W_k w}^2,
\]
dove $W_k = \operatorname{diag}([W_k]_{11}, \ldots, [W_k]_{nn})$ con $[W_k]_{ii} = |w_{k,i}|^{-1/2}$.

\begin{intubox}
  Questo è un'istanza del framework \textbf{Majorization-Minimization (MM)}:
  ad ogni passo si costruisce una funzione surrogate $\hat{f}_k(w)$ che
  \emph{maggiora} la funzione originale e la \emph{tocca} nel punto corrente $w_k$.
  Minimizzare la surrogate è più semplice e garantisce la discesa monotona della funzione originale.
\end{intubox}

\subsection{Sottoproblema di Minimi Quadrati Pesato}

Sostituendo l'approssimazione nell'obiettivo, il sottoproblema all'iterazione $k$ è:
\[
  w_{k+1} = \argmin_w \norm{Xw - y}^2 + 2\lambda\norm{W_k w}^2.
\]
Espandendo e ponendo il gradiente a zero:

\begin{keybox}[Sistema Lineare dell'IRLS]
  \[
    \underbrace{(A + 2\lambda W_k\T W_k)}_{=:\, Q_k}\, w_{k+1} = b,
    \qquad A = X\T X,\quad b = X\T y.
  \]
  La matrice $Q_k$ è \textbf{SPD} per ogni $k$ (se $\lambda > 0$), quindi si risolve con
  \textbf{fattorizzazione di Cholesky}.
\end{keybox}

\subsection{Aggiornamento della Matrice Peso}

\[
  [W_k]_{ii} = \frac{1}{\sqrt{\max(|w_{k,i}|,\;\varepsilon_{\text{thr}})}}, \qquad i = 1, \ldots, n,
\]
dove $\varepsilon_{\text{thr}} > 0$ è una soglia per la stabilità numerica (evita divisioni per zero).

\begin{intubox}
  Componenti $|w_{k,i}|$ grandi $\Rightarrow$ $[W_k]_{ii}$ piccolo $\Rightarrow$ poca penalizzazione
  $\Rightarrow$ il componente rimane grande.
  Componenti $|w_{k,i}|$ piccoli $\Rightarrow$ $[W_k]_{ii}$ grande $\Rightarrow$ forte penalizzazione
  $\Rightarrow$ il componente viene spinto verso zero.
  Questo implementa la \emph{sparsità adattiva} tipica della norma $\ell_1$.
\end{intubox}

\subsection{Algoritmo Completo (A1)}

\begin{algobox}[Algoritmo A1: IRLS]
  \textbf{Pre-computazione:} $A \leftarrow X\T X$,\quad $b \leftarrow X\T y$
  \vspace{0.5em}
  \begin{pseudocode}
    \pcline{$w_0 \leftarrow$ soluzione di $Aw = b$ (minimi quadrati puri, es.\ con Cholesky)}
    \pcline{\textbf{for} $k = 0, 1, 2, \ldots, k_{\max}$:}
    \pcline{\pcindent \textbf{for} $i = 1, \ldots, n$:}
    \pcline{\pcindentII $[W_k]_{ii} \leftarrow \max(|w_{k,i}|,\;\varepsilon_{\text{thr}})^{-1/2}$
      \hfill \textit{aggiorna pesi}}
    \pcline{\pcindent \textbf{end for}}
    \pcline{\pcindent $Q_k \leftarrow A + 2\lambda W_k\T W_k$
      \hfill \textit{matrice SPD}}
    \pcline{\pcindent Risolvi $Q_k\, w_{k+1} = b$ con Cholesky
      \hfill \textit{sottoproblema LS pesato}}
    \pcline{\pcindent \textbf{if} criterio di arresto soddisfatto: \textbf{break}}
    \pcline{\textbf{end for}}
    \pcline{\textbf{return} $w_{k+1}$}
  \end{pseudocode}
\end{algobox}

\subsection{Convergenza}

\begin{proposition}
  Per $\lambda > 0$ e $X$ a rango pieno, l'IRLS converge. La convergenza è
  \textbf{lineare (Q-lineare)} in un intorno della soluzione $w_*$:
  \[
    \norm{w_{k+1} - w_*} \leq r\, \norm{w_k - w_*}, \quad r < 1.
  \]
  Il tasso $r$ dipende dalla struttura della soluzione (grado di sparsità di $w_*$).
\end{proposition}

\begin{remark}
  Il costo per iterazione è $O(n^2)$ per aggiornare $W_k$ e $Q_k$, più $O(n^3/3)$ per
  la fattorizzazione di Cholesky. Nella pratica, il numero di iterazioni esterne è piccolo
  (convergenza rapida vicino alla soluzione) e IRLS risulta molto efficiente.
\end{remark}

% ================================================================
\section{Riepilogo: Connessioni tra gli Argomenti}
\label{sec:riepilogo}
% ================================================================

Questo manuale copre gli argomenti necessari per implementare e comprendere il Progetto 25,
organizzati nel seguente schema di dipendenze:

\begin{center}
  \renewcommand{\arraystretch}{1.2}
  \begin{tabular}{lll}
    \toprule
    \textbf{Argomento} & \textbf{Sezione} & \textbf{Uso nel Progetto} \\
    \midrule
    Minimi Quadrati & \ref{sec:ls} & Formulazione base ELM, sottoproblemi IRLS \\
    Equazioni Normali & \ref{sec:ls} & $A\T Aw = A\T y$, struttura di $Q_k$ in A1 \\
    Cholesky & \ref{sec:chol} & Solver interno di A1 per ogni iterazione \\
    Gradiente Coniugato & \ref{sec:cg} & Alternativa iterativa per sistemi sparsi \\
    Convessità / $L$-lisciezza & \ref{sec:opt} & Garanzie di ottimalità globale \\
    Subdifferenziale & \ref{sec:opt} & Condizioni di ottimalità per LASSO \\
    Complessità non liscia & \ref{sec:nonsmooth} & Perché servono metodi speciali \\
    Passo di Polyak & \ref{sec:nonsmooth} & Stepsize ottimale in A2 \\
    Target Level & \ref{sec:a2} & Stima adattiva di $f_*$ in A2 \\
    Direzione Deflessa & \ref{sec:a2} & Direzione migliorata in A2 \\
    MM Framework & \ref{sec:a1} & Giustificazione teorica di IRLS \\
    \bottomrule
  \end{tabular}
\end{center}

\vspace{1em}

\begin{keybox}[Percorso di Studio Consigliato]
  Sezione~\ref{sec:elm} (ELM + LASSO)
  $\;\to\;$ Sezione~\ref{sec:ls} (minimi quadrati)
  $\;\to\;$ Sezione~\ref{sec:chol} (Cholesky)
  $\;\to\;$ Sezione~\ref{sec:cg} (CG)
  $\;\to\;$ Sezione~\ref{sec:opt} (convessità)
  $\;\to\;$ Sezione~\ref{sec:nonsmooth} (ottimizzazione non liscia)
  $\;\to\;$ Sezione~\ref{sec:a1} (IRLS)
  $\;\to\;$ Sezione~\ref{sec:a2} (Sottogradiente Deflesso).
\end{keybox}

\end{document}
